% =====================================================================
%  template.tex  --  scheletro per un nuovo argomento
%
%  Come usarlo:
%    1. crea una cartella per l'argomento (es. DINAMICA/)
%    2. copia lì questo file e rinominalo (es. dinamica.tex)
%    3. sostituisci i testi segnaposto e compila da quella cartella
%
%  Stile e palette stanno nella cartella superiore (la root della repo):
%    ../appunti-base.sty   ../palette-fisica.tex
%  Compilazione: pdfLaTeX  (latexmk -pdf nomefile.tex)
%  Guida completa ad ambienti e comandi: STILE.md
% =====================================================================
\documentclass[11pt]{report}

\usepackage{../appunti-base}
\input{../palette-fisica}        % sempre DOPO appunti-base

% ---------- Salto pagina prima dei titoli ----------
% Di default una sezione va a pagina nuova se in fondo alla pagina restano
% meno di ~12 righe. Togli il commento qui sotto per disattivarlo in tutto
% il documento: le sezioni proseguono sulla stessa pagina.
% \nonewpage
% (nel testo: \nonewpage disattiva da lì in poi, \sinewpage riattiva)

% ---------- Dati della copertina ----------
\universita{Università di Pisa}
\corso{Corso di Laurea in Fisica}
\autore{Cristiano Battini}
\professore{Prof. Fabrizio Cei}
% \aggiornato{8 ottobre 2026}    % se omesso usa la data di compilazione

\begin{document}

\copertina{Titolo dell'argomento}
\tableofcontents

% =====================================================================
\chapter{Titolo del capitolo}
% =====================================================================

\lezione{GG mese AAAA}

\section{Titolo della sezione}

Testo degli appunti. Un \chiave{concetto chiave}, una \evid{parte evidenziata},
un passaggio logico \implica conseguenza.

\begin{definizione}{Nome}
  Testo della definizione.
\end{definizione}

\begin{teorema}{Nome della legge}
  Enunciato.
\end{teorema}

\begin{formula}
  x(t) = x_0 + v\,t
\end{formula}

\begin{dimostrazione}
  Passaggi.
\end{dimostrazione}

\begin{osservazione}
  Testo dell'osservazione.
\end{osservazione}

\subsection{Titolo della sottosezione}

\begin{esempio}{Titolo}
  Testo dell'esempio.
\end{esempio}

\begin{nota}
  Testo dell'NB.
\end{nota}

\begin{delucidazione}
  Nota a margine o scritta rossa degli appunti.
\end{delucidazione}

% ---------- Schema a frecce ----------
\begin{schema}
  \node[nodo]                 (a) {primo};
  \node[nodoevid, right=of a] (b) {secondo};
  \draw[freccia] (a) -- (b);
\end{schema}

% ---------- Disegno ----------
\begin{disegno}
  \begin{tikzpicture}[disegno]
    \draw[asse] (0,0) -- (4,0) node[right] {$t$};
    \draw[asse] (0,0) -- (0,2.5) node[above] {$x$};
    \draw[curva] (0,0.4) -- (3.5,2.1);
  \end{tikzpicture}
\end{disegno}

% ---------- Parti da ricontrollare ----------
% \illeggibile            \illeggibile[ipotesi]            \incerto{testo dubbio}

% =====================================================================
%  Riepilogo di fine blocco (facoltativo, sempre a pagina nuova)
% =====================================================================
% \riepilogo{Titolo del blocco}
%
% \begin{mappa}
%   \node[nodoevid]          (m) {Concetto};
%   \node[nodo, below=of m]  (c) {Caso};
%   \draw[freccia] (m) -- (c);
% \end{mappa}
%
% \begin{formulario}
%   x(t) = x_0 + v\,t & Cosa dice fisicamente \\
% \end{formulario}

\end{document}
